\scanpage{88}
\mspage{43}
under the following condition

(123) (Simp) For $\forall\ G\in\mathcal M$, the ordered set of quotients \inserted{$G'$} of $G$ is \deleted{\unsure{isomorphic} to a \unsure{simplex}}\inserted{a ``simplex'' i.e. isomorphic} to the set of subsets of a finite set.

If $A_G$ denotes the set of quotients of $G$ which are $\ne G$ and which are \deleted{\unsure{minimal|maximal}}\inserted{\unsure{maximal|minimal}} among these, then the \inserted{ordered} set of quotients of $G$ is therefore isomorphic to the set of subsets of $A_G$. \deleted{If} We denote by $G/C$ the quotient of $G$ associated with $A_G\setminus C$ (the smaller the larger $C$ is)\unsure{.}

\margin{\uline{NB} It would be necessary to show how this axiomatization applies to $\forall$ situation of multiplicities, \deleted{\unsure{topological or étale|ordinary topological}}\inserted{relative analytic} with a normal crossings divisor.... (except that $S\Pi M_G$ is an extension of $\Pi M_G$ not by a twisted $\mathbb Z^C$, but by a quotient of such a one.\unsure{-|..}}

We assume given for $\forall\ G\in\operatorname{Ob}\mathcal G$ two groupoids $S\Pi M_G$ and $\Pi M_G$, and a morphism
\[
\tag{124}
\diagram{\begin{tikzcd}[ampersand replacement=\&,column sep=1em]
S\Pi M_G \& \overset{\varphi_G}{\unsure{\rightharpoondown|{-}}} \& \Pi M_G.
\end{tikzcd}}
\]
We assume that $\varphi_G$ makes $S\Pi M_G$ an extension of $\Pi M_G$ by $\mathbb Z^{A(G)}$. (in fact, this is an additional structure on $S\Pi M_G$). We assume that $G\mapsto\Pi M_G$ is a functor in $G$ \uuline{for isomorphisms} (not for arbitrary morphisms)
\margin{87}
